\documentclass[../../main.tex]{subfiles}
\graphicspath{{\subfix{../../image/}}}

\begin{document}

\section*{习题五}\label{section:ch05-习题五}
\addcontentsline{toc}{section}{习题五}

\begin{exercise}\label{exercise:ch05-01}
$\xi$ 为非负随机变量，若 $\boldsymbol{E}\mathrm{e}^{a\xi}<\infty$（$a>0$），则对任意 $x>0$，
\begin{align*}
P\{\xi\geqslant x\}\leqslant\mathrm{e}^{-ax}\boldsymbol{E}\mathrm{e}^{a\xi}.
\end{align*}
\end{exercise}

\begin{exercise}\label{exercise:ch05-02}
若 $h(x)\geqslant0$，$\xi$ 为随机变量，且 $\boldsymbol{E}h(\xi)<\infty$，则关于任何 $C>0$，
\begin{align*}
P\{h(\xi)\geqslant C\}\leqslant C^{-1}\boldsymbol{E}h(\xi).
\end{align*}
\end{exercise}

\begin{exercise}\label{exercise:ch05-03}
（单边 Chebyshev 不等式）设 $\xi$ 为随机变量，$\boldsymbol{E}\xi=0$，$\boldsymbol{D}\xi=\sigma^{2}<\infty$，则对任何一个 $a>0$，试证
\begin{align*}
P\{\xi\geqslant a\}\leqslant\frac{\sigma^{2}}{\sigma^{2}+a^{2}}.
\end{align*}
\end{exercise}

\begin{exercise}\label{exercise:ch05-04}
设 $\{\xi_n,n\geqslant1\}$ 是独立随机变量序列，对所有 $n\geqslant1$，$\boldsymbol{D}\xi_n$ 存在且 $\dfrac{\boldsymbol{D}\xi_n}{n}\to0$，试证：$\{\xi_n,n\geqslant1\}$ 服从大数定律。
\end{exercise}

\begin{exercise}\label{exercise:ch05-05}
若 $\xi_k$ 的分布列为
\begin{center}
\begin{tabular}{c|cc}
$\xi_k$ & $\sqrt{\ln k}$ & $-\sqrt{\ln k}$\\ \hline
$P$ & $\dfrac12$ & $\dfrac12$
\end{tabular}
\end{center}
试证大数定律适用于独立随机变量序列 $\{\xi_k\}$。
\end{exercise}

\begin{exercise}\label{exercise:ch05-06}
验证概率分布如下给定的独立随机变量序列是否满足 Markov 条件：
\begin{enumerate}
\item $P\{X_k=\pm2^{k}\}=\dfrac12$；
\item $P\{X_k=\pm2^{k}\}=2^{-(2k+1)}$，$P\{X_k=0\}=1-2^{-2k}$；
\item $P\{X_k=\pm k\}=\dfrac12k^{-\frac12}$，$P\{X_k=0\}=1-k^{-\frac12}$。
\end{enumerate}
\end{exercise}

\begin{exercise}\label{exercise:ch05-07}
若 $\xi_k$ 具有有限方差，服从同一分布，但各 $k$ 间，$\xi_k$ 和 $\xi_{k+1}$ 有相关，而 $\xi_k,\xi_l$（$|k-l|\geqslant2$）是独立的，证明这时对 $\{\xi_k\}$ 大数定律成立。
\end{exercise}

\begin{exercise}\label{exercise:ch05-08}
（Bernstein 定理）已知随机变量序列 $\xi_1,\xi_2,\cdots$ 的方差有界：$\boldsymbol{D}\xi_n\leqslant C$，并且当 $|i-j|\to\infty$ 时，相关系数 $r_{ij}\to0$，证明对 $\{\xi_n\}$ 成立大数定律。
\end{exercise}

\begin{exercise}\label{exercise:ch05-09}
（Khinchin 定理）对随机变量序列 $\{\xi_i\}$，若记 $\eta_n=\dfrac1n(\xi_1+\cdots+\xi_n)$，$a_n=\dfrac1n(\boldsymbol{E}\xi_1+\cdots+\boldsymbol{E}\xi_n)$，则 $\{\xi_i\}$ 服从大数定律的充要条件是
\begin{align*}
\lim_{n\to\infty}\boldsymbol{E}\left\{\frac{(\eta_n-a_n)^{2}}{1+(\eta_n-a_n)^{2}}\right\}=0.
\end{align*}
\end{exercise}

\begin{exercise}\label{exercise:ch05-10}
用 Stirling 公式证明：当 $n\to\infty$，$m\to\infty$，$n-m\to\infty$，而 $\dfrac{m}{n}\to0$ 时，
\begin{align*}
\binom{2n}{n-m}\left(\frac12\right)^{2n}\sim\frac{1}{\sqrt{\pi n}}\mathrm{e}^{-m^{2}/n}.
\end{align*}
\end{exercise}

\begin{exercise}\label{exercise:ch05-11}
用正态逼近计算 $b(5;500,0.01)$ 及 $b(40;10000,0.005)$ 并与精确值比较。
\end{exercise}

\begin{exercise}\label{exercise:ch05-12}
某计算机系统有 $120$ 个终端，每个终端有 $5\%$ 时间在使用，若各个终端使用与否是相互独立的，试求有 $10$ 个或更多终端在使用的概率。
\end{exercise}

\begin{exercise}\label{exercise:ch05-13}
求证，在 $x>0$ 时，有不等式
\begin{align*}
\frac{x}{1+x^{2}}\mathrm{e}^{-x^{2}/2}\leqslant\int_{x}^{\infty}\mathrm{e}^{-t^{2}/2}\,\mathrm{d}t\leqslant\frac1x\mathrm{e}^{-x^{2}/2}.
\end{align*}
\end{exercise}

\begin{exercise}\label{exercise:ch05-14}
用 De Moivre--Laplace 定理证明，在 Bernoulli 试验中，若 $0<p<1$，则不管 $K$ 是如何大的常数，总有
\begin{align*}
P\{|\mu_n-np|<K\}\to0\qquad(n\to\infty).
\end{align*}
\end{exercise}

\begin{exercise}\label{exercise:ch05-15}
用 Chebyshev 不等式确定当掷一均匀铜币时，需投多少次才能保证使得正面出现的频率在 $0.4$ 至 $0.6$ 之间的概率不小于 $90\%$，并用正态逼近计算同一问题。
\end{exercise}

\begin{exercise}\label{exercise:ch05-16}
用 Chebyshev 不等式及 De Moivre--Laplace 极限定理估计下面概率：
\begin{align*}
P\!\left\{\left|\frac{\mu_n}{n}-p\right|\geqslant\varepsilon\right\}
\end{align*}
并进行比较。这里 $\mu_n$ 是 $n$ 次 Bernoulli 试验中成功总次数，$p$ 为每次成功的概率。
\end{exercise}

\begin{exercise}\label{exercise:ch05-17}
现有一大批种子，其中良种占 $1/6$，今在其中任选 $6000$ 粒，试问在这些种子中良种所占的比例与 $1/6$ 之差小于 $1\%$ 的概率是多少？
\end{exercise}

\begin{exercise}\label{exercise:ch05-18}
种子中良种占 $1/6$，我们有 $99\%$ 的把握断定在 $6000$ 粒种子中良种所占的比例与 $1/6$ 之差是多少？这时相应的良种粒数落在哪个范围内。
\end{exercise}

\begin{exercise}\label{exercise:ch05-19}
若飞机乘客购票后按期搭机的概率为 $p$，各乘客的行动假定是独立的，试问一架 $200$ 座飞机售出 $202$ 张机票不发生超座的概率。对 $p=0.97,0.96,0.95$，计算上述概率。
\end{exercise}

\begin{exercise}\label{exercise:ch05-20}
设分布函数列 $\{F_n(x)\}$ 弱收敛于连续的分布函数 $F(x)$，试证这收敛对 $x\in\mathbf{R}^{1}$ 是一致的。
\end{exercise}

\begin{exercise}\label{exercise:ch05-21}
设 $\{F_n(x)\}$ 为一列正态分布函数，收敛于分布函数 $F(x)$，试证 $F(x)$ 也是正态分布函数。
\end{exercise}

\begin{exercise}\label{exercise:ch05-22}
试证若正态随机变量序列依概率收敛，则其数学期望与方差也收敛。
\end{exercise}

\begin{exercise}\label{exercise:ch05-23}
若 $\bm{X}_n$ 为多维正态随机向量，$\bm{X}_n\xrightarrow{P}\bm{X}$，试证 $\bm{X}$ 为正态向量。
\end{exercise}

\begin{exercise}\label{exercise:ch05-24}
若 $X_n$ 的概率分布为
\begin{center}
\begin{tabular}{c|cc}
$X_n$ & $0$ & $n$\\ \hline
$P$ & $1-\dfrac1n$ & $\dfrac1n$
\end{tabular}
\end{center}
试证相应的分布函数列收敛，但矩不收敛。
\end{exercise}

\begin{exercise}\label{exercise:ch05-25}
（Slutsky）随机变量序列 $\{\xi_n\}$ 具有分布函数列 $\{F_n(x)\}$，且 $F_n(x)\to F(x)$，又 $\{\eta_n\}$ 依概率收敛于常数 $C>0$，试证：(1) $\zeta_n=\xi_n+\eta_n$ 的分布函数收敛于 $F(x-C)$；(2) $\zeta_n=\dfrac{\xi_n}{\eta_n}$ 的分布函数收敛于 $F(Cx)$。
\end{exercise}

\begin{exercise}\label{exercise:ch05-26}
试证：
\begin{enumerate}
\item $X_n\xrightarrow{P}X\Rightarrow X_n-X\xrightarrow{P}0$；
\item $X_n\xrightarrow{P}X$，$X_n\xrightarrow{P}Y\Rightarrow P\{X=Y\}=1$；
\item $X_n\xrightarrow{P}X\Rightarrow X_n-X_m\xrightarrow{P}0$（$n,m\to\infty$）；
\item $X_n\xrightarrow{P}X$，$Y_n\xrightarrow{P}Y\Rightarrow X_n\pm Y_n\xrightarrow{P}X\pm Y$；
\item $X_n\xrightarrow{P}X$，$k$ 是常数 $\Rightarrow kX_n\xrightarrow{P}kX$；
\item $X_n\xrightarrow{P}X\Rightarrow X_n^{2}\xrightarrow{P}X^{2}$；
\item $X_n\xrightarrow{P}a$，$Y_n\xrightarrow{P}b$，$a,b$ 是常数 $\Rightarrow X_nY_n\xrightarrow{P}ab$；
\item $X_n\xrightarrow{P}1\Rightarrow X_n^{-1}\xrightarrow{P}1$；
\item $X_n\xrightarrow{P}a$，$Y_n\xrightarrow{P}b$，$a,b$ 是常数，$b\neq0\Rightarrow X_nY_n^{-1}\xrightarrow{P}ab^{-1}$；
\item $X_n\xrightarrow{P}X$，$Y$ 是随机变量 $\Rightarrow X_nY\xrightarrow{P}XY$；
\item $X_n\xrightarrow{P}X$，$Y_n\xrightarrow{P}Y\Rightarrow X_nY_n\xrightarrow{P}XY$。
\end{enumerate}
\end{exercise}

\begin{exercise}\label{exercise:ch05-27}
设 $X_n\xrightarrow{P}X$，而 $g$ 是 $\mathbf{R}^{1}$ 上的连续函数，试证 $g(X_n)\xrightarrow{P}g(X)$。
\end{exercise}

\begin{exercise}\label{exercise:ch05-28}
若 $\{X_n\}$ 是单调下降的正随机变量序列，且 $X_n\xrightarrow{P}0$，试证 $X_n\xrightarrow{\mathrm{a.s.}}0$。
\end{exercise}

\begin{exercise}\label{exercise:ch05-29}
若 $X_1,X_2,\cdots$ 是独立随机变量序列，其特征函数均为 $\varphi(t)$，$\mu$ 是整值随机变量，$P\{\mu=k\}=p_k$，且与 $\{X_i\}$ 独立，求 $\eta=X_1+X_2+\cdots+X_{\mu}$ 的特征函数。
\end{exercise}

\begin{exercise}\label{exercise:ch05-30}
若 $f(t)$ 是非负定函数，试证：(1) $f(0)$ 是实的，且 $f(0)\geqslant0$；(2) $f(-t)=\overline{f(t)}$；(3) $|f(t)|\leqslant f(0)$。
\end{exercise}

\begin{exercise}\label{exercise:ch05-31}
某理发店为每个顾客的服务时间服从均值为 $\dfrac13$（小时）的指数分布，可认为对每个顾客的服务是相互独立的。
\begin{enumerate}
\item 求为对 $100$ 个顾客服务，总共需要 $31$ 小时至 $35$ 小时的概率；
\item 以 $95\%$ 的概率在 $32$ 小时之内可服务完几个顾客？
\item 找 $\Delta$，使该店对 $100$ 个顾客的服务时间在 $(33.33-\Delta,33.33+\Delta)$ 之间的概率大于 $95\%$。
\end{enumerate}
\end{exercise}

\begin{exercise}\label{exercise:ch05-32}
若总体 $\xi$ 的数学期望 $\boldsymbol{E}\xi=m$，$\boldsymbol{D}\xi=\sigma^{2}$，抽容量为 $n$ 的样本，求其平均值 $\bar{\xi}$，为使 $P\{|\bar{\xi}-m|<0.1\sigma\}\geqslant95\%$，问 $n$ 应取多大值？
\end{exercise}

\begin{exercise}\label{exercise:ch05-33}
用特征函数法直接证明 De Moivre--Laplace 积分极限定理。
\end{exercise}

\begin{exercise}\label{exercise:ch05-34}
若 $\{\xi_n,n=1,2,\cdots\}$ 为相互独立随机变量序列，具有相同分布
\begin{align*}
P\{\xi_n=1\}=\frac12,\quad P\{\xi_n=0\}=\frac12,
\end{align*}
而 $\eta_n=\sum_{k=1}^{n}\dfrac{\xi_k}{2^{k}}$，试证 $\eta_n$ 的分布收敛于 $[0,1]$ 上的均匀分布。
\end{exercise}

\begin{exercise}\label{exercise:ch05-35}
用特征函数法证明二项分布的 Poisson 逼近定理。
\end{exercise}

\begin{exercise}\label{exercise:ch05-36}
用特征函数法证明，Poisson 分布当 $\lambda\to\infty$ 时渐近正态分布。
\end{exercise}

\begin{exercise}\label{exercise:ch05-37}
若 $\{X_i\}$ 是独立同分布随机变量序列，其分布分别为：(1) $[-a,a]$ 上均匀分布；(2) Poisson 分布；(3) $\Gamma$ 分布。记
\begin{align*}
Y_n=\frac{\sum_{i=1}^{n}(X_i-\boldsymbol{E}X_i)}{\sqrt{\sum_{i=1}^{n}\boldsymbol{D}X_i}},
\end{align*}
试计算 $Y_n$ 的特征函数，并求 $n\to\infty$ 时的极限。
\end{exercise}

\begin{exercise}\label{exercise:ch05-38}
设 $\{X_n\}$ 独立同分布，$P\{X_n=2^{k-2\ln k}\}=2^{-k}$（$k=1,2,\cdots$），则大数定律成立。
\end{exercise}

\begin{exercise}\label{exercise:ch05-39}
若 $\{X_i\}$ 是相互独立的随机变量序列，均服从 $N(0,1)$，试证
\begin{align*}
W_n=\sqrt{n}\,\frac{X_1+\cdots+X_n}{\sqrt{X_1^{2}+\cdots+X_n^{2}}}\quad\text{及}\quad U_n=\frac{X_1+\cdots+X_n}{\sqrt{X_1^{2}+\cdots+X_n^{2}}}
\end{align*}
渐近正态分布 $N(0,1)$。
\end{exercise}

\begin{exercise}\label{exercise:ch05-40}
设 $X_1,X_2,\cdots$ 是独立随机变量序列，均服从 $[0,1]$ 均匀分布，令
\begin{align*}
Z_n=\left(\prod_{i=1}^{n}X_i\right)^{1/n},
\end{align*}
试证 $Z_n\xrightarrow{P}C$，这里 $C$ 是常数，并求 $C$。
\end{exercise}

\begin{exercise}\label{exercise:ch05-41}
若 $\{X_i\}$ 是独立同分布随机变量序列，$\boldsymbol{E}X_i=m$，而 $f(x)$ 是一个有界的连续函数，试证
\begin{align*}
\lim_{n\to\infty}\boldsymbol{E}\!\left[f\!\left(\frac{X_1+\cdots+X_n}{n}\right)\right]=f(m).
\end{align*}
\end{exercise}

\begin{exercise}\label{exercise:ch05-42}
若 $\{X_i\}$ 是独立同分布、具有有限二阶矩的随机变量序列，试证
\begin{align*}
\frac{2}{n(n+1)}\cdot\sum_{i=1}^{n}iX_i\xrightarrow{P}\boldsymbol{E}X_1.
\end{align*}
\end{exercise}

\begin{exercise}\label{exercise:ch05-43}
设 $X_1,X_2,\cdots$ 相互独立，均服从 Cauchy 分布 $p(x)=\dfrac1\pi\cdot\dfrac{1}{1+x^{2}}$，试证它们不满足 Khinchin 关于大数定律的充要条件（见本章习题 9），即要指出，当 $n\to\infty$ 时
\begin{align*}
\boldsymbol{E}\!\left[\frac{\left(\sum_{i=1}^{n}X_i\right)^{2}}{n^{2}+\left(\sum_{i=1}^{n}X_i\right)^{2}}\right]\not\to0.
\end{align*}
\end{exercise}

\begin{exercise}\label{exercise:ch05-44}
（Weierstrass 定理的概率论证明）设 $f(x)$ 是 $[0,1]$ 上连续函数，利用概率论方法证明：必存在多项式序列 $\{B_n(x)\}$，在 $[0,1]$ 上一致收敛于 $f(x)$。（提示：定义 Bernstein 多项式
\begin{align*}
B_n(x)=\sum_{m=0}^{n}\binom{n}{m}x^{m}(1-x)^{n-m}f\!\left(\frac{m}{n}\right)
\end{align*}
并利用大数定律。）
\end{exercise}

\begin{exercise}\label{exercise:ch05-45}
设 $\{X_n\}$ 是独立随机变量序列，试证 $X_n\xrightarrow{\mathrm{a.s.}}0$ 的充要条件为对任意 $\varepsilon>0$，有 $\sum_{n=1}^{\infty}P\{|X_n|\geqslant\varepsilon\}<\infty$。
\end{exercise}

\begin{exercise}\label{exercise:ch05-46}
试证独立同分布随机变量序列，若存在有限的四阶中心矩，则强大数定律成立。
\end{exercise}

\begin{exercise}\label{exercise:ch05-47}
本章习题 6 的各个独立随机变量序列是否满足强大数定律？
\end{exercise}

\begin{exercise}\label{exercise:ch05-48}
举例说明 Borel--Cantelli 引理 (i) 之逆不成立。
\end{exercise}

\begin{exercise}\label{exercise:ch05-49}
设 $\{X_n\}$ 是相互独立且具有有限方差的随机变量序列，若
\begin{align*}
\sum_{n=1}^{\infty}\frac{\boldsymbol{D}X_n}{n^{2}}<\infty,
\end{align*}
则必有
\begin{align*}
\lim_{n\to\infty}\frac{1}{n^{2}}\sum_{k=1}^{n}\boldsymbol{D}X_k=0.
\end{align*}
\end{exercise}

\begin{exercise}\label{exercise:ch05-50}
设 $f(x)$ 和 $g(x)$ 在闭区间 $[0,1]$ 上连续，且满足 $0\leqslant f(x)<Cg(x)$，这里 $C$ 是一个正常数，则成立
\begin{align*}
\lim_{n\to\infty}\int_{0}^{1}\cdots\int_{0}^{1}\frac{f(x_1)+f(x_2)+\cdots+f(x_n)}{g(x_1)+g(x_2)+\cdots+g(x_n)}\,\mathrm{d}x_1\mathrm{d}x_2\cdots\mathrm{d}x_n
=\frac{\int_{0}^{1}f(x)\,\mathrm{d}x}{\int_{0}^{1}g(x)\,\mathrm{d}x}.
\end{align*}
\end{exercise}

\begin{exercise}\label{exercise:ch05-51}
设 $X_1,X_2,\cdots$ 是独立随机变量序列，对它成立中心极限定理，则对 $\{X_n\}$ 成立大数定律的充要条件为 $\boldsymbol{D}(X_1+\cdots+X_n)=o(n^{2})$。
\end{exercise}

\begin{exercise}\label{exercise:ch05-52}
设 $X_1,X_2,\cdots$ 是独立同分布随机变量序列，且 $\dfrac{\sum_{k=1}^{n}X_k}{\sqrt{n}}$ 对每一个 $n=1,2,\cdots$ 有相同分布，那么，若 $\boldsymbol{E}X_i=0$，$\boldsymbol{D}X_i=1$，则 $X_i$ 必须是 $N(0,1)$ 变量。
\end{exercise}

\begin{exercise}\label{exercise:ch05-53}
设 $\{X_k\}$ 是独立随机变量序列，且 $X_k$ 服从 $N(0,2^{-k})$，试证序列 $\{X_k\}$：(1) 成立中心极限定理；(2) 不满足 Feller 条件；(3) 不满足 Lindeberg 条件，从而说明 Lindeberg 条件并不是中心极限定理成立的必要条件。
\end{exercise}

\begin{exercise}\label{exercise:ch05-54}
若 $\{X_k\}$ 是独立随机变量序列，$X_1$ 服从 $[-1,1]$ 均匀分布，对 $k=2,3,\cdots$，$X_k$ 服从 $N(0,2^{k-1})$，证明对 $\{X_k\}$ 成立中心极限定理但不满足 Feller 条件。
\end{exercise}

\begin{exercise}\label{exercise:ch05-55}
在 Poisson 试验中，第 $i$ 次试验时事件 $A$ 出现的概率为 $p_i$，不出现的概率为 $q_i$，各次试验是独立的，以 $\nu_n$ 记前 $n$ 次试验中事件 $A$ 出现的次数。试证：
\begin{enumerate}
\item $\dfrac{\nu_n-\boldsymbol{E}\nu_n}{n}\xrightarrow{P}0$；
\item 对 $\{\nu_n\}$ 成立中心极限定理的充要条件是 $\sum_{i=1}^{\infty}p_iq_i=+\infty$。
\end{enumerate}
\end{exercise}

\begin{exercise}\label{exercise:ch05-56}
设 $\{X_k\}$ 是独立随机变量序列，$X_k$ 服从 $[-k,k]$ 均匀分布，问对 $\{X_k\}$ 能否用中心极限定理？
\end{exercise}

\begin{exercise}\label{exercise:ch05-57}
试问对下列独立随机变量序列，Lyapunov 定理是否成立？
\begin{enumerate}
\item $X_k$：
\begin{center}
\begin{tabular}{c|cc}
 & $-\sqrt{k}$ & $\sqrt{k}$\\ \hline
$P$ & $\dfrac12$ & $\dfrac12$
\end{tabular}
\end{center}
\item $X_k$：
\begin{center}
\begin{tabular}{c|ccc}
 & $-k^{\alpha}$ & $0$ & $k^{\alpha}$\\ \hline
$P$ & $\dfrac13$ & $\dfrac13$ & $\dfrac13$
\end{tabular}
\end{center}
$\alpha>0$。
\end{enumerate}
\end{exercise}

\begin{exercise}\label{exercise:ch05-58}
求证：当 $n\to\infty$ 时
\begin{align*}
\frac{\left(\dfrac n2\right)^{\frac n2}}{\Gamma\!\left(\dfrac n2\right)}\int_{0}^{1+\sqrt{\frac2n}}z^{\frac n2-1}\mathrm{e}^{-\frac{nz}{2}}\,\mathrm{d}z\to\frac{1}{\sqrt{2\pi}}\int_{-\infty}^{t}\mathrm{e}^{-z^{2}/2}\,\mathrm{d}z.
\end{align*}
\end{exercise}

\begin{exercise}\label{exercise:ch05-59}
独立随机变量序列 $\{\xi_i\}$，对一切 $k$，$\xi_k$ 以概率 $\dfrac12$ 分别取值 $\pm k^{s}$，
\begin{enumerate}
\item 试证当 $s<\dfrac12$ 时，大数定律成立；
\item 试利用中心极限定理证明：当 $s\geqslant\dfrac12$ 时，大数定律不成立。
\end{enumerate}
\end{exercise}

\begin{exercise}\label{exercise:ch05-60}
用概率论方法证明：当 $n\to\infty$ 时，$\mathrm{e}^{-n}\sum_{k=0}^{n}\dfrac{n^{k}}{k!}\to\dfrac12$。
\end{exercise}

\end{document}
